\section{Theory of Operation}

\subsection{Overview}
The \textbf{Pcs} core implements the Physical Coding Sublayer (PCS) for 10GBASE-R Ethernet as defined in IEEE 802.3 Clause 49. It serves as the bridge between the Media Access Control (MAC) layer’s XGMII interface and the Physical Medium Attachment (PMA) / SERDES interface.

The core performs 64b/66b line coding, which includes data encapsulation, scrambling for DC balance, and block synchronization. The architecture is fully duplex, featuring independent Transmit (TX) and Receive (RX) paths.
The PCS performs the following functions:
\begin{itemize}
    \item 64b/66b block encoding and decoding
    \item Scrambling and descrambling for DC balance
    \item Block synchronization and alignment (Block Lock)
    \item High Bit Error Rate (BER) monitoring
    \item XGMII control character mapping
    \item Hardware PRBS31 test pattern support
\end{itemize}

\subsection{Architecture}
\begin{figure}[H]
\centering
\begin{tikzpicture}[
    block/.style={rectangle, draw=space-blue, fill=fog-grey, thick, minimum width=4cm, minimum height=1cm, align=center, rounded corners=2pt},
    arrow/.style={-{Stealth[scale=1.2]}, thick, space-blue}
]

% Center Nodes
\node (mac) [text=space-blue, font=\bfseries] at (0, 0) {MAC (XGMII Interface)};
\node (serdes) [text=space-blue, font=\bfseries] at (0, -4.5) {SERDES (Physical PHY)};

% TX Path (Left)
\node (txEnc) [block] at (-3, -2.25) {\textbf{TX Path} \\ (Encoder/Scrambler)};

% RX Path (Right)
\node (rxDec) [block] at (3, -2.25) {\textbf{RX Path} \\ (Sync/Descrambler)};

% Arrows
\draw [arrow] (mac.west) -| (txEnc.north);
\draw [arrow] (txEnc.south) |- (serdes.west);
\draw [arrow] (serdes.east) -| (rxDec.south);
\draw [arrow] (rxDec.north) |- (mac.east);

\end{tikzpicture}
\caption{PCS Top-Level Architecture}\label{fig:pcs_arch}
\end{figure}



\subsection{Transmit Path Logic}
The TX datapath converts 64-bit XGMII data and 8-bit control characters into a scrambled 66-bit bitstream.

\textbf{Transmit Pipeline:}
\begin{figure}[H]
\centering
\begin{tikzpicture}[
    node distance=0.8cm,
    block/.style={rectangle, draw=crockadile-green, fill=fog-grey, thick, minimum width=4.5cm, minimum height=0.8cm, align=center},
    arrow/.style={-{Stealth[scale=1.2]}, thick, crockadile-green}
]

\node (xgmii) [text=crockadile-green, font=\bfseries] {XGMII TX};
\node (enc) [block, below=0.6cm of xgmii] {XGMII Encoder};
\node (scram) [block, below=of enc] {Scrambler};
\node (format) [block, below=of scram] {64b/66b Formatter};
\node (serdes) [text=crockadile-green, font=\bfseries, below=0.6cm of format] {SERDES TX};

\draw [arrow] (xgmii) -- (enc);
\draw [arrow] (enc) -- (scram);
\draw [arrow] (scram) -- (format);
\draw [arrow] (format) -- (serdes);

\end{tikzpicture}
\caption{Transmit Datapath Pipeline}\label{fig:pcs_tx_pipe}
\end{figure}

\subsubsection{64b/66b Encoding}
The \textit{XgmiiEncoder} compresses XGMII characters into one of 15 valid 64b/66b block types. Each block is prepended with a 2-bit synchronization header:
\begin{itemize}
    \item \textbf{\texttt{b10}:} Data Block (Contains only data payload).
    \item \textbf{\texttt{b01}:} Control Block (Contains start/terminate characters, ordered sets, or idles).
\end{itemize}
If the encoder receives an illegal XGMII sequence, it asserts \textit{txBadBlock} and substitutes the data with an Error block (\texttt{0xFE}).

\subsubsection{Scrambling}
To ensure sufficient transition density and DC balance for high-speed serial clock recovery, the payload is processed by a self-synchronizing scrambler. It utilizes the polynomial $G(x) = 1 + x^{39} + x^{58}$. Similar to the MAC's CRC implementation, the \textbf{Pcs} core calculates the scrambler matrix at elaboration time to perform 64-bit parallel scrambling in a single clock cycle.
The scrambler may be disabled through configuration.

\subsubsection{Gearbox Handshaking}
Because the 64b/66b encoding process increases the data rate by a factor of $66/64$, the core implements gearbox handshaking. The \textit{txGbxReqStall} signal provides backpressure to the MAC, ensuring that the TX path maintains the correct timing for the extra header bits without overflowing internal buffers.

\subsection{Receive Path Logic}
The Receive path recovers the 64b/66b block boundaries from the raw SERDES input and reverses the encoding process.

\textbf{Receive Pipeline:}
\begin{figure}[H]
\centering
\begin{tikzpicture}[
    node distance=0.8cm,
    block/.style={rectangle, draw=orange, fill=fog-grey, thick, minimum width=4.5cm, minimum height=0.8cm, align=center},
    arrow/.style={-{Stealth[scale=1.2]}, thick, orange}
]

\node (serdes) [text=orange, font=\bfseries] {SERDES RX};
\node (sync) [block, below=0.6cm of serdes] {Block Synchronizer};
\node (descram) [block, below=of sync] {Descrambler};
\node (dec) [block, below=of descram] {XGMII Decoder};
\node (xgmii) [text=orange, font=\bfseries, below=0.6cm of dec] {XGMII RX};

\draw [arrow] (serdes) -- (sync);
\draw [arrow] (sync) -- (descram);
\draw [arrow] (descram) -- (dec);
\draw [arrow] (dec) -- (xgmii);

\end{tikzpicture}
\caption{Receive Datapath Pipeline}\label{fig:pcs_rx_pipe}
\end{figure}

\subsubsection{Block Synchronization and Bitslip}
Upon power-up or link loss, the \textit{PcsRxFrameSync} module enters a search state. It monitors the 2-bit sync headers for valid patterns (\texttt{01} or \texttt{10}). If invalid headers are detected, the core asserts \textit{serdesRxBitslip}, instructing the SERDES to shift the bit alignment. \textit{rxBlockLock} is asserted once 64 consecutive valid headers are identified.

\subsubsection{Descrambling and Decoding}
Once block lock is achieved, the \textit{PcsRxInterface} utilizes a self-synchronizing descrambler to recover the original payload. The \textit{XgmiiDecoder} then expands the 64b/66b blocks back into standard XGMII data and control characters for the MAC.

\subsection{Link Monitoring and Diagnostics}
The PCS continuously monitors the bitstream for protocol violations. The primary error indicators are summarized in Table~\ref{table:pcs_errors}.

\subsubsection{Bit Error Rate (BER) Monitor}
The \textit{PcsRxBerMon} module continuously validates sync headers within a moving 125$\mu$s window. If more than 16 errors are detected in this window, \textit{rxHighBer} is asserted to notify the system of a degraded physical link.

\subsubsection{Watchdog Timer}
The \textit{PcsRxWatchdog} acts as a protocol-level supervisor. It monitors for sustained "Bad Block" or "Sequence Error" counts. If the decoder reports excessive errors despite having block lock, the watchdog can trigger a \textit{serdesRxResetReq} to force a full link re-alignment.

\subsubsection{PRBS31 Support}
For hardware validation, the core includes a Pseudo-Random Binary Sequence (PRBS31) generator and checker. When \textit{cfgTxPrbs31Enable} is high, the TX path transmits a $2^{31}-1$ pattern. The RX path can independently check this pattern and report bit errors via the 7-bit \textit{rxErrorCount} register.

\begin{table}[H]
\centering
\begin{tabular}{|l|l|}
\hline
\textbf{Signal} & \textbf{Description} \\ \hline \hline
\texttt{rxBadBlock} & Asserted when an invalid 64b/66b block type is detected. \\ \hline
\texttt{rxSequenceError} & Indicates a violation in the XGMII control character sequence. \\ \hline
\texttt{rxHighBer} & Asserted when the Bit Error Rate exceeds the $10^{-4}$ threshold. \\ \hline
\end{tabular}
\caption{PCS Error Indicators}\label{table:pcs_errors}
\end{table}

\subsection{System Integration}
To integrate the \textbf{Pcs} core into a larger system, the following interfaces must be connected:
\begin{enumerate}
    \item \textbf{MAC Interface:} Connects to the MAC core via the 64-bit XGMII interface.
    \item \textbf{SERDES Interface:} Connects to high-speed FPGA transceivers (GTH/GTY) via a 64-bit (or 32-bit) raw data interface.
    \item \textbf{Handshaking:} Uses gearbox stall signals to regulate data flow, ensuring that the additional 2 bits of 64b/66b overhead do not cause buffer overflows in the upstream MAC.
\end{enumerate}

Figure~\ref{fig:mac_pcs_combined} illustrates how the \textbf{Pcs} core fits within a complete 10GBASE-R architecture, showing the flow of data from the application layer down to the physical medium and back up.

\begin{figure}[H]
\centering
\begin{tikzpicture}[
    block/.style={rectangle, draw=space-blue, fill=fog-grey, thick, minimum width=4cm, minimum height=1cm, align=center, rounded corners=2pt},
    arrow/.style={-{Stealth[scale=1.2]}, thick, space-blue},
    interface/.style={text=space-blue, font=\bfseries}
]

% Center Nodes (The Boundaries)
\node (app) [interface] at (1, 0) {Application Logic};
\node (xgmii) [interface] at (-6, -4.5) {XGMII Interface};
\node (serdes) [interface] at (1, -9) {SERDES (Physical PHY)};

% Left Column (TX Path - Flowing Down)
\node (axiTx) at (-2.2, -1.2) {AXI-Stream TX};
\node (macTx) [block] at (-2.2, -2.5) {\textbf{MAC TX} \\ (Axis2Xgmii64)};
\node (xgmiiTx) at (-2.2, -4.5) {64-bit Data + 8-bit Ctrl};
\node (pcsTx) [block] at (-2.2, -6.5) {\textbf{PCS TX} \\ (Encoder / Scrambler)};
\node (phyTx) at (-2.2, -7.8) {64b/66b TX};

% Right Column (RX Path - Flowing Up)
\node (axiRx) at (4.2, -1.2) {AXI-Stream RX};
\node (macRx) [block] at (4.2, -2.5) {\textbf{MAC RX} \\ (Xgmii2Axis64)};
\node (xgmiiRx) at (4.2, -4.5) {64-bit Data + 8-bit Ctrl};
\node (pcsRx) [block] at (4.2, -6.5) {\textbf{PCS RX} \\ (Sync / Descrambler)};
\node (phyRx) at (4.2, -7.8) {64b/66b RX};

% --- Arrows for TX Path (Flowing Down) ---
\draw [arrow] (app.west) -| (axiTx.north);
\draw [arrow] (axiTx.south) -- (macTx.north);
\draw [arrow] (macTx.south) -- (xgmiiTx.north);
\draw [arrow] (xgmiiTx.south) -- (pcsTx.north);
\draw [arrow] (pcsTx.south) -- (phyTx.north);
\draw [arrow] (phyTx.south) |- (serdes.west);

% --- Arrows for RX Path (Flowing Up) ---
\draw [arrow] (serdes.east) -| (phyRx.south);
\draw [arrow] (phyRx.north) -- (pcsRx.south);
\draw [arrow] (pcsRx.north) -- (xgmiiRx.south);
\draw [arrow] (xgmiiRx.north) -- (macRx.south);
\draw [arrow] (macRx.north) -- (axiRx.south);
\draw [arrow] (axiRx.north) |- (app.east);

% --- Optional Domain Separators ---
% These subtle dashed lines visually separate the MAC layer from the PCS layer
\draw [dashed, space-blue!40, thick] (-8, -3.75) -- (8, -3.75);
\draw [dashed, space-blue!40, thick] (-8, -5.25) -- (8, -5.25);

\end{tikzpicture}
\caption{Combined MAC and PCS 10GBASE-R Architecture}\label{fig:mac_pcs_combined}
\end{figure}